%%% ABSTRACT

\todo[inline]{Do I need to add references to the abstract here? Needs consistency across the thesis!}

Systems memory consolidation involves the transfer of memories across brain regions and the transformation of memory content. For example, declarative memories that transiently depend on the hippocampal formation are transformed into long-term memory traces in neo-cortical networks, and procedural memories are transformed within cortico-striatal networks. These consolidation processes are thought to rely on replay and repetition of recently acquired memories, but the cellular and network mechanisms that mediate the changes of memories are poorly understood. Here, we suggest that systems memory consolidation could arise from Hebbian plasticity in networks with parallel synaptic pathways --- two ubiquitous features of neural circuits in the brain. We explore this hypothesis in mathematical analyses and computational models to illustrate how memories are transferred across circuits and how this could mediate the semantization of episodic memories. These modelling results are in quantitative agreement with lesion studies in rodents. A hierarchical iteration of the mechanism yields power-law forgetting -- as observed in psychophysical studies in humans, and it yields gradients of episodic and semantic memory strength across the hierarchy. Moreover, the results suggest that memory replay is only necessary for the transfer of episodic but not semantic memories. The predicted circuit mechanism thus bridges spatial scales from single cells to cortical areas, time scales from milliseconds to years, and connects the transfer of memories to their transformation into semantic knowledge.

\todo[inline]{Removed the linear approximation sentence}
